\documentclass[UTF8]{ctexart}
\usepackage[a4paper,margin=2.6cm]{geometry}
\begin{document}

\section*{第三部分：最短路与其他图论概念}
本文默认图 $G$ 是一个有向图，并指定一个源点 $s$，而且所有顶点都能从 $s$ 到达。若无特别说明，图中不含重边和自环。无向图可以看成有向图：每条无向边替换成两条方向相反、长度相同的弧。本文关于最短路问题的结论对无向图同样成立。

设 $n$ 和 $m$ 分别表示图的顶点数和弧数。为简化一些复杂度表述，作者假设 $n>2$；在连通图里这意味着 $n=O(m)$。

每条弧 $vw$ 都有一个非负实长度 $c(vw)$。一条路径的长度就是其各条弧长度之和。若从 $v$ 到 $w$ 的所有路径中，某条路径长度最小，就称它是最短路径。由于弧长非负，只要两点之间存在路径，就一定存在一条不重复顶点的最短简单路。记 $d^*(v)$ 为从源点 $s$ 到 $v$ 的最短路长度，也就是 $v$ 的真实距离。以 $s$ 为根、所有根到点路径都是最短路的生成树，称为最短路树。

研究最短路算法时，还会用到一些额外概念。若一个有向无环图存在一个顶点全序，使得每条弧都从较前的点指向较后的点，这个全序就是拓扑序。对一般的有向图 $G$ 和源点 $s$，若存在一种给弧赋非负长度的方法，使得所有真实距离都两两不同，并且某个顶点全序恰好按真实距离严格递增，那么这个全序称为距离顺序。作者用列表 $L$ 表示这样的顺序，并记图 $G$（带源点 $s$）的距离顺序总数为 $D$。

作者给出一个很有用的判定：一个顶点全序是距离顺序，当且仅当除了源点 $s$ 之外的每个顶点，都至少有一条来自前驱顶点的入弧。换句话说，只要每个非源点都能从更靠前的点接到，它就是某种距离顺序。

在给定一个顶点顺序后，如果某条弧 $vw$ 满足 $v$ 排在 $w$ 前面，就称它是这个顺序的前向弧。于是，距离顺序也可以理解成：每个非源点都至少有一条来自前方的入弧。作者记前向弧数最多的距离顺序所具有的前向弧数为 $F$。任何距离顺序至少会有 $n-1$ 条前向弧，因为其中至少包含一棵生成树的边。对无向图而言，每条无向边对应两条相反方向的弧，其中恰好一条是前向弧，所以这时 $F=m$。

如果在某组具体弧长下，某个距离顺序还与这些弧长对应的真实距离一致地非递减，那么它就是一个真实距离顺序。本文主要研究的问题，就是如何计算图 $G$ 的真实距离顺序。

作者这样定义真实距离顺序，是为了允许弧长为零。Dijkstra 算法对任意非负弧长都正确，所以本文希望包含这种一般情况。如果不要求每个非源点都有一条前向入弧，那么在所有弧长都为零的图里，任意顶点顺序都可以看成非递减顺序，这显然不是作者想要的概念。

最后，作者引入的是从源点出发路径上的“不可避免顶点”概念，主要包括支配点和瓶颈点。

对两个不同顶点 $v$ 和 $w$，如果从 $s$ 到 $w$ 的每条路径都经过 $v$，就说 $v$ 支配 $w$。$w$ 的直接支配点是：在所有支配点里，那个被其他支配点共同支配的唯一顶点。每个顶点都有直接支配点，这些关系形成一棵支配树。支配树包含图中的所有顶点，但它的边未必都在原图中。一个顶点的所有支配点，正好是它在支配树中的真祖先。

若弧 $vw$ 满足 $w$ 支配 $v$，则称它是无用弧；否则称为有用弧。无用弧不可能成为任何距离顺序中的前向弧，也不可能出现在从 $s$ 到任意顶点的简单最短路上，所以求最短路时可以忽略它们。无用弧可以通过先求支配树，再标记那些指向祖先的弧来找到，总时间是 $O(m)$。

不过，Dijkstra 算法实际上并不需要显式找出这些无用弧，也不需要显式计算支配树。本文的比较最优算法只依赖一个更弱的概念：瓶颈点。定义顶点 $v$ 的层数 $\ell(v)$ 为从 $s$ 到 $v$ 的路径上最少顶点数。于是对任意弧 $vw$，总有 $\ell(v)+1\ge \ell(w)$。如果某个顶点 $v$ 是它所在层里唯一的顶点，就称它是瓶颈点。瓶颈点会支配所有更高层的顶点，而且它们在支配树上落在同一条路径上。层数和瓶颈点都可以通过一次从 $s$ 出发的广度优先搜索在 $O(m)$ 时间内求出。

\end{document}
